% Summary document for supervisor (concise) -- Overleaf-ready
\documentclass[12pt]{article}
\usepackage{amsmath, amssymb, graphicx, bm, geometry}
\geometry{margin=1in}

\title{\textbf{Summary of PhD Framework and Initial Experimental Plan}}
\author{PhD Candidate: [Your Name] \\ Supervisor: [Supervisor Name] \\ Institution: [Institution]}
\date{\today}

\begin{document}
%\maketitle

\section{Overview of the PhD Direction}
This PhD aims to develop a unified and practical framework for \textbf{real-time aerodynamic state estimation and prediction} for bluff bodies, ranging from canonical geometries to simplified human-body postures. Classical CFD provides high-fidelity flow fields, but it is too computationally expensive for real-time applications and cannot be directly integrated with sparse experimental sensors.

The thesis therefore proposes a \textbf{reduced-order aerodynamic state model}, informed by coherent structures, equipped with:
\begin{itemize}
    \item a low-dimensional \textbf{state representation},
    \item a \textbf{predictive model} of state evolution,
    \item a \textbf{measurement operator} linking state to sparse sensor outputs, and
    \item a \textbf{corrector} based on data assimilation to ensure physical consistency.
\end{itemize}

This architecture remains the same across increasingly complex cases; only the specific modal basis, model parameters, and sensor configurations change.

\section{Research Goal}
To determine whether aerodynamic forces---especially drag---can be reconstructed and predicted in real time using only a small number of pressure sensors, and to establish a modelling strategy that generalises from simple to complex geometries.

\section{Three-Stage Progressive Validation}
The work proceeds through three validation stages:

\subsection*{1. Cylinder (Proof of Concept)}
A cylinder facing a wind wall is instrumented with four pressure sensors (front, rear, and two lateral positions). Measurements are collected at three flow velocities.

Objectives:
\begin{itemize}
    \item Determine whether sparse pressure measurements contain enough information to reconstruct the drag coefficient.
    \item Build and compare:
    \begin{itemize}
        \item a direct regression model (baseline),
        \item a reduced-order model (modal or latent representation),
        \item a data-driven \textbf{state map} via clustering and transition probabilities.
    \end{itemize}
    \item Evaluate robustness across the three velocities (e.g., train on two, test on the third).
    \item Identify whether coherent aerodynamic ``states'' emerge from the pressure/drag signals.
\end{itemize}

Criteria for success:
\begin{itemize}
    \item Accurate reconstruction of drag (target: mean error $<5$--$10\%$).
    \item Consistent pressure-to-drag mapping across velocities.
    \item Interpretable states or modes representing dominant flow structures.
\end{itemize}

If the method fails at this stage, the approach must be revised before proceeding.

\subsection*{2. Simplified 3D Dummy}
A simplified human-like geometry is instrumented with sensors at selected locations (e.g., hands, knees, chest). Only a restricted set of posture variations is considered.

Goals:
\begin{itemize}
    \item Extend the modelling framework beyond canonical flows.
    \item Characterise how posture changes affect aerodynamic state and forces.
    \item Validate the ability of the predictor--corrector system to track changes in posture.
\end{itemize}

\subsection*{3. Human Body with Limited Degrees of Freedom}
The final stage focuses on realistic aerodynamic behaviour of a human subject with a small number of controlled posture changes.

Goals:
\begin{itemize}
    \item Build a reduced-order model for a low-dimensional posture space.
    \item Predict drag in real time.
    \item Demonstrate the ability to suggest transitions toward lower-drag states.
\end{itemize}

\section{Mathematical Foundation (Concise)}
The high-dimensional flow field $\mathbf{q}(\mathbf{x}, t)$ is approximated using a reduced state:
\[
  \mathbf{q}(t) \approx \bar{\mathbf{q}} + \sum_{i=1}^r x_i(t)\,\boldsymbol{\phi}_i,
\]
where $\boldsymbol{\phi}_i$ are POD/SPOD modes or latent vectors.

A reduced-order dynamic model governs state evolution:
\[
  \mathbf{x}_{k+1} = f(\mathbf{x}_k; \theta) + \boldsymbol{\xi}_k.
\]

Pressure sensors provide measurements through:
\[
  \mathbf{y}_k = h(\mathbf{x}_k; \phi) + \boldsymbol{\eta}_k.
\]

Drag is estimated from the reduced state via a mapping $g$:
\[
  \hat{C}_D(k) = g(\mathbf{x}_k; \psi).
\]

Data assimilation (EKF/EnKF) combines predictions and pressure data to form a real-time estimator.

\section{Initial Experimental Plan: Cylinder Test}
\subsection*{Experimental Setup}
\begin{itemize}
    \item Wind wall with a fixed circular cylinder.
    \item Three flow velocities: $U_1$, $U_2$, $U_3$.
    \item Four pressure sensors located at:
    \begin{itemize}
        \item front stagnation region,
        \item rear base region,
        \item left and right sides.
    \end{itemize}
    \item Force balance to measure drag $C_D(t)$.
\end{itemize}

\subsection*{What We Aim to Validate}
\begin{enumerate}
    \item Whether pressure signals alone can reconstruct the drag coefficient reliably.
    \item Whether a low-dimensional latent or modal state emerges naturally from the data.
    \item Whether clustering reveals meaningful aerodynamic states and transitions.
    \item Whether models trained at certain velocities generalise to others.
    \item Whether the predictor--corrector formulation stabilises drag estimation.
\end{enumerate}

\subsection*{Success Criteria}
\begin{itemize}
    \item Drag prediction error below $5$--$10\%$.
    \item Stable performance across all three velocities.
    \item Identification of coherent aerodynamic signatures.
    \item Clear evidence that the modelling architecture is viable beyond canonical flows.
\end{itemize}

\end{document}
